\documentclass[12pt,a4paper]{article}

% ── Packages ────────────────────────────────────────────────────────────────
\usepackage[utf8]{inputenc}
\usepackage[T1]{fontenc}
\usepackage{lmodern}
\usepackage[margin=2.5cm, headheight=15pt]{geometry}
\usepackage{amsmath, amssymb, amsthm}
\usepackage{mathtools}
\usepackage{enumitem}
\usepackage{booktabs}
\usepackage{multicol}
\usepackage{xcolor}
\usepackage{titlesec}
\usepackage{fancyhdr}
\usepackage{hyperref}

% ── Colours & formatting ────────────────────────────────────────────────────
\definecolor{primary}{RGB}{0,70,127}
\definecolor{secondary}{RGB}{180,20,20}
\definecolor{shadecolor}{RGB}{240,245,255}

\hypersetup{colorlinks=true, linkcolor=primary, urlcolor=primary}

\titleformat{\section}[block]
  {\large\bfseries\color{primary}}
  {\thesection.}{0.5em}{}[\titlerule]

\titleformat{\subsection}[runin]
  {\normalsize\bfseries\color{secondary}}
  {\thesubsection.}{0.4em}{}[.\quad]

% ── Page style ───────────────────────────────────────────────────────────────
\pagestyle{fancy}
\fancyhf{}
\lhead{\textcolor{primary}{\textbf{MTH253 -- Calculus III}}}
\rhead{\textcolor{primary}{Exercise Sheet}}
\cfoot{\thepage}
\renewcommand{\headrulewidth}{0.4pt}

% ── Custom environments ──────────────────────────────────────────────────────
\newtheoremstyle{exercise}
  {8pt}{6pt}{\normalfont}{0pt}{\bfseries\color{primary}}{.}{0.5em}{}
\theoremstyle{exercise}
\newtheorem{exercise}{Exercise}[section]

\newtheoremstyle{hint}
  {4pt}{4pt}{\itshape}{0pt}{\bfseries\color{secondary}}{:}{0.4em}{}
\theoremstyle{hint}
\newtheorem*{hint}{Hint}

% ── Handy macros ─────────────────────────────────────────────────────────────
\newcommand{\R}{\mathbb{R}}
\newcommand{\dd}{\,\mathrm{d}}
\newcommand{\pd}[2]{\dfrac{\partial #1}{\partial #2}}
\newcommand{\pdd}[2]{\dfrac{\partial^{2} #1}{\partial #2^{2}}}
\newcommand{\grad}{\nabla}
\newcommand{\norm}[1]{\left\lVert #1 \right\rVert}
\newcommand{\abs}[1]{\left\lvert #1 \right\rvert}

% ────────────────────────────────────────────────────────────────────────────
\begin{document}

\begin{center}
  {\LARGE\bfseries\color{primary} Calculus III --- Exercise Sheet}\\[4pt]
  {\large\color{secondary} Chapter 1: Differentiation \quad|\quad
   Chapter 2: Multiple Integrals}\\[4pt]
  {\normalsize University of Science, VNU-HCM \hfill
   Academic Year 2024--2025}
\end{center}
\vspace{-4pt}
\hrule height 1.2pt
\vspace{12pt}

\tableofcontents
\newpage

% ============================================================
\section{Vectors and Geometry of Space}
% ============================================================

\begin{exercise}
Find the vector $\overrightarrow{AB}$ and its magnitude for each of the
following pairs of points.
\begin{enumerate}[label=(\alph*)]
  \item $A(1,-2,4)$, $B(3,0,-1)$.
  \item $A(-3,5,0)$, $B(2,-1,6)$.
\end{enumerate}
\end{exercise}

\begin{exercise}
Let $\mathbf{u}=\langle 2,-1,3\rangle$ and
$\mathbf{v}=\langle -1,4,2\rangle$.
\begin{enumerate}[label=(\alph*)]
  \item Compute the dot product $\mathbf{u}\cdot\mathbf{v}$.
  \item Find the angle $\theta$ between $\mathbf{u}$ and $\mathbf{v}$.
  \item Compute the cross product $\mathbf{u}\times\mathbf{v}$ and verify
        that it is orthogonal to both $\mathbf{u}$ and $\mathbf{v}$.
  \item Find a unit vector in the direction of $\mathbf{u}$.
\end{enumerate}
\end{exercise}

\begin{exercise}
\begin{enumerate}[label=(\alph*)]
  \item Find parametric and symmetric equations of the line passing
        through $P(-1,2,3)$ and parallel to the vector
        $\mathbf{v}=\langle 3,-1,5\rangle$.
  \item Find parametric equations of the line passing through the points
        $A(1,0,-2)$ and $B(4,3,1)$.
  \item Find the point where the line $x=1+2t,\;y=-t,\;z=3+t$ meets
        the plane $2x-y+z=7$.
\end{enumerate}
\end{exercise}

\begin{exercise}
\begin{enumerate}[label=(\alph*)]
  \item Find an equation of the plane passing through $P(2,-1,3)$ with
        normal vector $\mathbf{n}=\langle 1,5,-2\rangle$.
  \item Find an equation of the plane through the three points
        $P(0,1,1)$, $Q(1,0,1)$, and $R(1,1,0)$.
  \item Determine whether the planes $2x-3y+z=4$ and
        $4x-6y+2z=1$ are parallel, identical, or neither.
\end{enumerate}
\end{exercise}

\begin{exercise}
Describe and sketch the following surfaces in $\R^3$.
\begin{enumerate}[label=(\alph*)]
  \item $x^{2}+z^{2}=9$.
  \item $z=4-x^{2}-y^{2}$.
  \item $\dfrac{x^{2}}{4}+\dfrac{y^{2}}{9}+\dfrac{z^{2}}{16}=1$.
\end{enumerate}
\end{exercise}

% ============================================================
\section{Functions of Several Variables}
% ============================================================

\begin{exercise}
For each function, (i) find and sketch the domain, (ii) determine the
range, and (iii) describe three representative level curves.
\begin{enumerate}[label=(\alph*)]
  \item $f(x,y)=\sqrt{16-x^{2}-y^{2}}$.
  \item $f(x,y)=\ln(x+y-1)$.
  \item $f(x,y)=\dfrac{x^{2}-y^{2}}{x^{2}+y^{2}}$.
\end{enumerate}
\end{exercise}

\begin{exercise}
Find the domain and the level surfaces of
\[
  f(x,y,z)=\sqrt{9-x^{2}-y^{2}-z^{2}}.
\]
\end{exercise}

\begin{exercise}
Sketch the level curves of $f(x,y)=x^{2}+4y^{2}$ for the values
$k=0,1,4,9$.
\end{exercise}

% ============================================================
\section{Limits and Continuity}
% ============================================================

\begin{exercise}
Evaluate each limit, or show that it does not exist.
\begin{enumerate}[label=(\alph*)]
  \item $\displaystyle\lim_{(x,y)\to(1,2)}\bigl(x^{3}y-2xy^{2}+5\bigr)$.
  \item $\displaystyle\lim_{(x,y)\to(0,0)}\frac{5x^{2}y}{x^{2}+y^{2}}$.
  \item $\displaystyle\lim_{(x,y)\to(0,0)}\frac{x^{2}-xy}{\sqrt{x}-\sqrt{y}}$
        (where defined).
  \item $\displaystyle\lim_{(x,y)\to(0,0)}\frac{xy^{2}}{x^{2}+y^{4}}$.
\end{enumerate}
\end{exercise}

\begin{exercise}
Use the Squeeze Theorem to prove that
\[
  \lim_{(x,y)\to(0,0)}\frac{x^{2}y^{2}}{x^{2}+y^{2}}=0.
\]
\end{exercise}

\begin{exercise}
Discuss the continuity of the function
\[
  f(x,y)=
  \begin{cases}
    \dfrac{2xy}{x^{2}+y^{2}} & \text{if }(x,y)\neq(0,0),\\[6pt]
    0 & \text{if }(x,y)=(0,0).
  \end{cases}
\]
\end{exercise}

% ============================================================
\section{Partial Derivatives and Higher Derivatives}
% ============================================================

\begin{exercise}
Find all first-order partial derivatives of each function.
\begin{enumerate}[label=(\alph*)]
  \item $f(x,y)=x^{4}y-3x^{2}y^{3}+7y$.
  \item $g(x,y)=e^{x}\sin(xy)$.
  \item $h(x,y,z)=xy^{2}z^{3}+\ln(x+2y+3z)$.
  \item $w(s,t)=s\arctan(t/s)$.
\end{enumerate}
\end{exercise}

\begin{exercise}
Find all second-order partial derivatives of
$f(x,y)=x^{3}y^{2}-2x^{2}y^{4}+5xy$, and verify Clairaut's theorem
($f_{xy}=f_{yx}$).
\end{exercise}

\begin{exercise}
A function $f(x,y)$ satisfies Laplace's equation if
$f_{xx}+f_{yy}=0$.  Verify that each of the following functions is
harmonic.
\begin{enumerate}[label=(\alph*)]
  \item $f(x,y)=x^{3}-3xy^{2}$.
  \item $f(x,y)=e^{x}\cos y$.
  \item $f(x,y)=\ln(x^{2}+y^{2})$.
\end{enumerate}
\end{exercise}

\begin{exercise}
Let $f(x,y)=x^{2}e^{y}+y\sin x$.  Calculate $f_{xxy}$ and $f_{xyy}$.
\end{exercise}

% ============================================================
\section{Tangent Planes and Linear Approximations}
% ============================================================

\begin{exercise}
\begin{enumerate}[label=(\alph*)]
  \item Find the equation of the tangent plane to the surface
        $z=x^{2}+2y^{2}-3x$ at the point $(1,1,0)$.
  \item Find the equation of the tangent plane to the surface
        $z=e^{x-y}$ at the point $(2,2,1)$.
\end{enumerate}
\end{exercise}

\begin{exercise}
Find the linearisation $L(x,y)$ of $f(x,y)=\sqrt{20-x^{2}-7y^{2}}$
at $(2,1)$, and use it to estimate $f(1.95,1.05)$.
\end{exercise}

\begin{exercise}
The dimensions of a closed rectangular box are measured to be
$80\,\text{cm}$, $60\,\text{cm}$, and $50\,\text{cm}$, with a possible
error of $\pm 0.2\,\text{cm}$ in each measurement.  Use differentials
to estimate the maximum error in the calculated volume of the box.
\end{exercise}

\begin{exercise}
Find the total differential $dz$ for each function.
\begin{enumerate}[label=(\alph*)]
  \item $z=x^{4}y+x\sqrt{y}$.
  \item $z=\arctan\!\left(\dfrac{x}{y}\right)$.
\end{enumerate}
\end{exercise}

% ============================================================
\section{The Chain Rule}
% ============================================================

\begin{exercise}
Use the Chain Rule to find $dz/dt$.
\begin{enumerate}[label=(\alph*)]
  \item $z=x^{2}y+3xy^{4}$,\quad $x=\sin t$,\quad $y=e^{t}$.
  \item $z=\sqrt{x^{2}+y^{2}}$,\quad $x=e^{t}\cos t$,\quad $y=e^{t}\sin t$.
\end{enumerate}
\end{exercise}

\begin{exercise}
Use the Chain Rule to find $\partial z/\partial s$ and
$\partial z/\partial t$.
\begin{enumerate}[label=(\alph*)]
  \item $z=e^{2x+3y}$,\quad $x=s/t$,\quad $y=t/s$.
  \item $z=x^{2}\tan y$,\quad $x=st^{2}$,\quad $y=s^{2}t$.
\end{enumerate}
\end{exercise}

\begin{exercise}[Implicit Differentiation]
\begin{enumerate}[label=(\alph*)]
  \item Find $dy/dx$ if $x^{4}+x^{2}y^{2}+y^{4}=3$.
  \item Find $\partial z/\partial x$ and $\partial z/\partial y$ if
        $x^{2}+2y^{2}+3z^{2}+xyz=6$.
\end{enumerate}
\end{exercise}

% ============================================================
\section{Directional Derivatives and the Gradient}
% ============================================================

\begin{exercise}
Find the directional derivative of each function at the given point in
the given direction.
\begin{enumerate}[label=(\alph*)]
  \item $f(x,y)=x^{2}y^{3}-2x$,\; at $(1,-1)$,\;
        direction angle $\theta=\pi/4$.
  \item $f(x,y,z)=x^{2}yz+xz^{2}$,\; at $(1,1,1)$,\;
        $\mathbf{v}=\langle 1,2,-1\rangle$.
\end{enumerate}
\end{exercise}

\begin{exercise}
\begin{enumerate}[label=(\alph*)]
  \item Find the gradient of $f(x,y)=\sin(x+2y)$ and evaluate it at
        $(\pi/6,0)$.
  \item Find the maximum rate of change of
        $f(x,y,z)=\sqrt{x^{2}+y^{2}+z^{2}}$ at $(1,2,2)$ and the
        direction in which it occurs.
\end{enumerate}
\end{exercise}

\begin{exercise}
The temperature at a point $(x,y)$ in the plane is given by
$T(x,y)=200\,e^{-(x^{2}+y^{2})/100}$.
A bug is located at $(3,4)$.
\begin{enumerate}[label=(\alph*)]
  \item In which direction should the bug move to warm up most quickly?
  \item How fast does the temperature change in that direction?
  \item If the bug moves in the direction $\mathbf{u}=\langle 4,3\rangle/5$,
        what is the rate of change of temperature?
\end{enumerate}
\end{exercise}

% ============================================================
\section{Extrema of Functions of Several Variables}
% ============================================================

\begin{exercise}
Find and classify all critical points (local max, local min, or saddle
point) of each function.
\begin{enumerate}[label=(\alph*)]
  \item $f(x,y)=x^{2}+xy+y^{2}-3x$.
  \item $f(x,y)=4xy-x^{4}-y^{4}$.
  \item $f(x,y)=e^{x}(x+y^{2}+2y)$.
\end{enumerate}
\end{exercise}

\begin{exercise}
Find the absolute maximum and minimum values of $f(x,y)=x^{2}-2xy+2y$
on the closed triangular region $D$ with vertices $(0,0)$, $(3,0)$, and
$(3,3)$.
\end{exercise}

\begin{exercise}
Find three positive numbers $a$, $b$, $c$ whose sum is $100$ and whose
product is as large as possible.
\end{exercise}

% ============================================================
\section{Lagrange Multipliers}
% ============================================================

\begin{exercise}
Use Lagrange multipliers to find the extreme values of each function
subject to the given constraint.
\begin{enumerate}[label=(\alph*)]
  \item Maximise $f(x,y)=xy$ subject to $x+2y=8$.
  \item Minimise $f(x,y,z)=x^{2}+y^{2}+z^{2}$ subject to $x+y+z=12$.
  \item Find the extreme values of $f(x,y)=x^{2}+2y^{2}$ on the circle
        $x^{2}+y^{2}=1$.
\end{enumerate}
\end{exercise}

\begin{exercise}
Use Lagrange multipliers to find the point on the plane $2x+y+z=5$
closest to the origin.
\end{exercise}

\begin{exercise}
A rectangular box without a lid is to be made from $12\,\text{m}^{2}$
of cardboard.  Use Lagrange multipliers to find the dimensions that
maximise the volume of the box.
\end{exercise}

% ============================================================
\section{Double Integrals over Rectangles}
% ============================================================

\begin{exercise}
Evaluate the following iterated integrals.
\begin{enumerate}[label=(\alph*)]
  \item $\displaystyle\int_{0}^{3}\int_{1}^{2}
          \bigl(x^{2}y-2xy^{3}\bigr)\dd y\dd x$.
  \item $\displaystyle\int_{0}^{\pi}\int_{0}^{1}
          y\sin(xy)\dd y\dd x$.
  \item $\displaystyle\int_{0}^{1}\int_{0}^{1}
          xe^{xy}\dd y\dd x$.
\end{enumerate}
\end{exercise}

\begin{exercise}
Evaluate $\displaystyle\iint_{R}(x+2y)\dd A$ where
$R=[0,2]\times[0,3]$ in two ways: integrating first with respect to $y$
then $x$, and vice versa.  Verify you get the same answer.
\end{exercise}

\begin{exercise}
Find the volume of the solid bounded above by the plane
$z=2+x+3y$ and below by the rectangle
$R=[-1,1]\times[0,2]$ in the $xy$-plane.
\end{exercise}

% ============================================================
\section{Double Integrals over General Regions}
% ============================================================

\begin{exercise}
Evaluate each double integral over the described region $D$.
\begin{enumerate}[label=(\alph*)]
  \item $\displaystyle\iint_{D}(2x-y)\dd A$, where $D$ is bounded by the
        curves $y=x^{2}$ and $y=x+2$.
  \item $\displaystyle\iint_{D}xy\dd A$, where $D$ is bounded by
        $y=x$ and $y=x^{2}$.
  \item $\displaystyle\iint_{D}x\dd A$, where $D$ is the triangle with
        vertices $(0,0)$, $(1,0)$, and $(0,1)$.
\end{enumerate}
\end{exercise}

\begin{exercise}
Reverse the order of integration and evaluate.
\begin{enumerate}[label=(\alph*)]
  \item $\displaystyle\int_{0}^{1}\int_{x}^{1}\sin(y^{2})\dd y\dd x$.
  \item $\displaystyle\int_{0}^{4}\int_{\sqrt{y}}^{2}
          e^{x^{3}}\dd x\dd y$.
\end{enumerate}
\end{exercise}

\begin{exercise}
Use a double integral to find the area enclosed by the curves
$y=3x-x^{2}$ and $y=x$.
\end{exercise}

% ============================================================
\section{Double Integrals in Polar Coordinates}
% ============================================================

\begin{exercise}
Convert to polar coordinates and evaluate.
\begin{enumerate}[label=(\alph*)]
  \item $\displaystyle\iint_{D}\sqrt{x^{2}+y^{2}}\dd A$,\quad
        $D=\{(x,y):x^{2}+y^{2}\leq 4,\; x\geq 0\}$.
  \item $\displaystyle\iint_{D}e^{-(x^{2}+y^{2})}\dd A$,\quad
        $D$ is the disk of radius $a$ centred at the origin.
  \item $\displaystyle\int_{0}^{2}\int_{0}^{\sqrt{4-x^{2}}}
          (x^{2}+y^{2})^{3/2}\dd y\dd x$.
\end{enumerate}
\end{exercise}

\begin{exercise}
Find the volume of the solid that lies under the paraboloid
$z=9-x^{2}-y^{2}$, above the $xy$-plane, and inside the cylinder
$x^{2}+y^{2}=4$.
\end{exercise}

\begin{exercise}
Use polar coordinates to find the area of the region inside the
cardioid $r=1+\cos\theta$ and outside the circle $r=1$.
\end{exercise}

% ============================================================
\section{Triple Integrals}
% ============================================================

\begin{exercise}
Evaluate the following triple integrals.
\begin{enumerate}[label=(\alph*)]
  \item $\displaystyle\int_{0}^{2}\int_{0}^{x}\int_{0}^{x+y}
          6xy\dd z\dd y\dd x$.
  \item $\displaystyle\iiint_{B}xe^{y+2z}\dd V$, where
        $B=[0,1]\times[0,1]\times[0,1]$.
  \item $\displaystyle\iiint_{E}y\dd V$, where $E$ is the solid bounded
        by $z=0$, $z=x+y+5$, $x=0$, $x=1$, $y=0$, and $y=1$.
\end{enumerate}
\end{exercise}

\begin{exercise}
\begin{enumerate}[label=(\alph*)]
  \item Evaluate $\displaystyle\iiint_{E}z\dd V$, where $E$ is the solid
        tetrahedron with vertices $(0,0,0)$, $(1,0,0)$, $(0,1,0)$, and
        $(0,0,1)$.
  \item Find the volume of the solid enclosed by the paraboloids
        $z=x^{2}+y^{2}$ and $z=36-3x^{2}-3y^{2}$.
\end{enumerate}
\end{exercise}

% ============================================================
\section{Triple Integrals in Cylindrical Coordinates}
% ============================================================

\begin{exercise}
Use cylindrical coordinates to evaluate the following integrals.
\begin{enumerate}[label=(\alph*)]
  \item $\displaystyle\iiint_{E}\bigl(x^{2}+y^{2}\bigr)\dd V$, where $E$
        is the solid inside the cylinder $x^{2}+y^{2}=4$ and between the
        planes $z=-3$ and $z=5$.
  \item $\displaystyle\iiint_{E}\sqrt{x^{2}+y^{2}}\dd V$, where $E$ is
        bounded above by $z=4$ and below by the paraboloid
        $z=x^{2}+y^{2}$.
  \item $\displaystyle\int_{-2}^{2}\int_{-\sqrt{4-x^{2}}}^{\sqrt{4-x^{2}}}
          \int_{\sqrt{x^{2}+y^{2}}}^{2}(x^{2}+y^{2})\dd z\dd y\dd x$.
\end{enumerate}
\end{exercise}

\begin{exercise}
A solid hemisphere of radius $a$ has its flat face in the $xy$-plane.
Using cylindrical coordinates, find the volume of the hemisphere.
\end{exercise}

% ============================================================
\section{Triple Integrals in Spherical Coordinates}
% ============================================================

\begin{exercise}
Use spherical coordinates to evaluate the following integrals.
\begin{enumerate}[label=(\alph*)]
  \item $\displaystyle\iiint_{B}e^{(x^{2}+y^{2}+z^{2})^{3/2}}\dd V$,
        where $B$ is the unit ball $x^{2}+y^{2}+z^{2}\leq 1$.
  \item $\displaystyle\iiint_{H}(x^{2}+y^{2}+z^{2})^{2}\dd V$, where
        $H$ is the upper hemisphere of radius $2$.
  \item $\displaystyle\iiint_{E}z\dd V$, where $E$ lies above the cone
        $z=\sqrt{x^{2}+y^{2}}$ and inside the sphere $x^{2}+y^{2}+z^{2}=4$.
\end{enumerate}
\end{exercise}

\begin{exercise}
\begin{enumerate}[label=(\alph*)]
  \item Find the volume of the solid that lies inside the sphere
        $x^{2}+y^{2}+z^{2}=4$ and above the cone $z=\sqrt{x^{2}+y^{2}}$.
  \item Use spherical coordinates to evaluate
        $\displaystyle\iiint_{E}\bigl(x^{2}+y^{2}\bigr)\dd V$, where $E$
        is the ball of radius $3$.
\end{enumerate}
\end{exercise}

% ============================================================
\section{Mixed Challenge Problems}
% ============================================================

\begin{exercise}
The temperature in a metal plate is given by
$T(x,y)=x^{2}-xy+y^{2}$.
\begin{enumerate}[label=(\alph*)]
  \item Find all critical points of $T$ and classify them.
  \item Find the maximum and minimum temperatures on the disk
        $x^{2}+y^{2}\leq 4$.
\end{enumerate}
\end{exercise}

\begin{exercise}
Let $f(x,y)=x^{2}+4xy+y^{3}$.
\begin{enumerate}[label=(\alph*)]
  \item Use implicit differentiation to find $dy/dx$ if
        $f(x,y)=0$.
  \item Find the gradient $\grad f$ at $(1,-1)$.
  \item Compute the directional derivative at $(1,-1)$ in the direction
        toward the origin.
\end{enumerate}
\end{exercise}

\begin{exercise}
\emph{(Gaussian Integral via double integrals)}  Using polar coordinates,
evaluate $I=\displaystyle\int_{-\infty}^{\infty}e^{-x^{2}}\dd x$ by
computing
\[
  I^{2}=
  \left(\int_{-\infty}^{\infty}e^{-x^{2}}\dd x\right)
  \!\left(\int_{-\infty}^{\infty}e^{-y^{2}}\dd y\right)
  =\iint_{\R^{2}}e^{-(x^{2}+y^{2})}\dd A.
\]
\end{exercise}

\begin{exercise}
Find the volume of the region $E$ bounded by
$z=x^{2}+y^{2}$ (paraboloid) and $x^{2}+y^{2}+z^{2}=6$
(sphere) using the most convenient coordinate system, justifying your
choice.
\end{exercise}

\begin{exercise}
A curve in the $xy$-plane is defined implicitly by
$F(x,y)=x\sin(y)-y\cos(x)=0$.
\begin{enumerate}[label=(\alph*)]
  \item Use the implicit function theorem to find $dy/dx$ near $(0,0)$.
  \item Write the equation of the tangent line to the curve at $(0,0)$.
\end{enumerate}
\end{exercise}

\bigskip
\begin{center}
  \rule{0.6\linewidth}{0.4pt}\\[4pt]
  {\small\itshape Good luck! Remember: mathematics is not a spectator sport.}
\end{center}

\end{document}
